\section{System design}\label{sec:design}

The pipeline turns a recording into a fixed size image and then classifies that image.
Three choices inside it are the subject of the study, and this section takes them in the
order the signal meets them: Section~\ref{sec:design-denoise} covers the wavelet
denoising applied to the waveform, Section~\ref{sec:design-spectro} covers the
spectrogram that is formed from the result, and Section~\ref{sec:design-nets} covers the
four networks the images are shown to. Everything else about the pipeline is held fixed
across conditions so that a difference in the outcome can be attributed to one of these
three.

\subsection{Wavelet denoising of the waveform}\label{sec:design-denoise}

A crowdsourced recording carries whatever noise the contributor's room and microphone
supplied. Wavelet shrinkage is an established way to suppress such noise while preserving
transients~\cite{donoho1994}, which matters here because a cough is a transient and a
method that smooths it away would destroy the thing being detected.

The waveform $x$ is decomposed to level $J = 6$ with the \texttt{sym4} wavelet using the
discrete wavelet transform, computed by the pyramidal algorithm of
Mallat~\cite{mallat1989}. This yields detail coefficients $d_{j,k}$ at scales
$j = 1, \dots, J$ and a residual approximation. Shrinkage is applied to the detail
coefficients only, and the signal is reconstructed from the shrunk coefficients.

Which shrinkage rule is used is not a detail that can be left implicit, because the
reported effect of denoising depends on it entirely. Three rules appear in this work and
they are distinguished throughout. The universal threshold of Donoho and
Johnstone~\cite{donoho1994}, applied with soft thresholding, sets one threshold for the
whole decomposition from an estimate $\hat{\sigma}$ of the noise level:
\begin{equation}
\lambda_{\mathrm{uni}} = \hat{\sigma}\sqrt{2\ln n},
\qquad
\eta_{\mathrm{soft}}(d) = \operatorname{sgn}(d)\,\max(|d| - \lambda, 0),
\label{eq:visu}
\end{equation}
where $n$ is the number of samples and $\hat{\sigma}$ is taken from the median absolute
deviation of the finest scale coefficients. The level adaptive rule of Chang, Yu and
Vetterli~\cite{chang2000} instead sets a threshold per scale from the ratio of the noise
variance to the estimated signal variance $\hat{\sigma}_{x,j}^{2}$ at that scale:
\begin{equation}
\lambda_{\mathrm{bayes},j} = \frac{\hat{\sigma}^{2}}{\hat{\sigma}_{x,j}},
\qquad
\hat{\sigma}_{x,j}^{2} = \max\!\left(\widehat{\operatorname{var}}(d_{j,\cdot})
  - \hat{\sigma}^{2},\, 0\right).
\label{eq:bayes}
\end{equation}
The third rule is the empirical Bayes procedure of Johnstone and
Silverman~\cite{johnstone2005}, which places a mixture prior on each coefficient,
\begin{equation}
d_{j,k} \sim (1 - w_j)\,\delta_0 + w_j\,\gamma,
\label{eq:eb}
\end{equation}
with $\delta_0$ a point mass at zero, $\gamma$ a heavy tailed density, and the mixing
weight $w_j$ estimated from the coefficients at that scale by marginal maximum
likelihood. The coefficient is then replaced by the median of its posterior.

This matters for interpreting the work being extended. Its implementation calls a
denoising routine without naming a rule, and the default of that routine is the empirical
Bayes procedure of Equation~\ref{eq:eb} with a Cauchy prior and the posterior median,
not the soft thresholding of Equation~\ref{eq:visu}. The denoising reported here
implements Equations~\ref{eq:visu} and~\ref{eq:bayes} directly and states which is in
use. It is therefore a replication under a specified denoiser rather than a reproduction
of that routine, and it is not expected to agree with it coefficient by coefficient. The
alternative, reporting a result under a rule that was never named, is what makes a
denoising comparison impossible to check.

Figure~\ref{fig:denoise} shows what the shrinkage removes, for one cough and one
non-cough recording. The panels showing the removed component are each normalised
independently, so they show where in time and frequency the shrinkage acted and not how
much energy it took: the magnitudes of two such panels are not comparable with one
another.

\subsection{Spectrogram formation and the frequency axis}\label{sec:design-spectro}

Each waveform is resampled to $16$~kHz and a short time Fourier transform is taken with a
$512$ point window and a hop of $256$ samples, giving $257$ frequency bins. Magnitudes
are converted to decibels. Two conditions differ only in what happens next, and this is
the second variable of the study.

In the linear condition the $257$ bins are used as they are. In the mel condition they are
mapped onto $128$ bands spaced on the mel scale of Stevens, Volkmann and
Newman~\cite{stevens1937}, which spaces bands according to perceived rather than physical
pitch and so allocates more resolution to low frequencies:
\begin{equation}
m(f) = 2595 \log_{10}\!\left(1 + \frac{f}{700}\right).
\label{eq:mel}
\end{equation}
The filterbank is triangular and follows the formulation used in hidden Markov model
toolkits. One property of it is worth stating because it bears on how much the mel
condition can differ from the linear one: at this configuration $128$ bands are requested
over $257$ bins, and one of the $128$ triangular filters contains no bin centre at all and
is therefore empty. The mel axis here is oversampled relative to the resolution the
transform provides, so the two conditions differ less in information content than the
change of axis suggests.

Both conditions produce a single channel image, which is resized to $224 \times 224$ and
replicated across three channels so that networks pretrained on photographs can accept
it. The conversion from decibels to pixel values is per recording, which means absolute
loudness is discarded and only the pattern within a recording survives.

\subsection{Networks}\label{sec:design-nets}

Four networks are used, chosen so that capacity and initialisation vary while the input
does not. Three are pretrained on ImageNet~\cite{russakovsky2015} and have their final
classification layer replaced by a two way layer: ResNet18~\cite{he2016}, which supplies
residual connections at modest depth, MobileNetV2~\cite{sandler2018}, which uses inverted
residuals and depthwise convolutions and is an order of magnitude smaller in multiply
accumulate operations, and DenseNet201~\cite{huang2017}, which is far deeper and
concatenates features across layers. The fourth is a compact network trained from random
initialisation, consisting of three blocks of convolution, batch normalisation,
rectification and max pooling, followed by a fully connected layer with dropout.

The fourth network exists to separate two things the work being extended could not. Its
baseline differed from its denoised condition in architecture, in pretraining, in the
frequency axis, and in the sample rate at once, so no single cause could be assigned to
the gap between them. Here the from scratch network sees exactly the images the
pretrained networks see, so the contribution of pretraining is isolated from the
contribution of the input.

All four are trained with Adam~\cite{kingma2015} in PyTorch~\cite{paszke2019}. The batch
size is the only hyperparameter that varies by network, and it varies for a hardware
reason rather than a modelling one: DenseNet201 activations at this input size do not fit
alongside the resident images in the available memory at a larger batch.
